% Discussion for item1_3 -- prose goes here.
% scripts/make_latex.py will NEVER overwrite this file.

The two builds do produce different results, as the table shows.

By default \texttt{gcc} generates a position-independent executable (PIE) whose
base address is randomised by ASLR on each execution, causing the address of
static \texttt{y} to change across runs. Compiling with \texttt{-no-pie} fixes
the program's load address, pinning \texttt{y} to \texttt{0x404018}. The values
(6, 7, 8) remain unaffected, as storage duration is governed by language
semantics, not by binary loading.

Notably, the non-static \texttt{y} moves across runs even under
\texttt{-no-pie}. Because stack layout is randomised independently by the
kernel, disabling PIE fixes only the text and data segments, not the stack.
